\documentclass[10pt,a4paper]{article}
\usepackage{fullpage}
\usepackage[utf8]{inputenc}
\usepackage[T1]{fontenc}
\usepackage{amsmath}
\usepackage{amssymb}
\usepackage{graphicx}
\usepackage[english]{babel}
\usepackage{url}
\usepackage{hyperref}
\usepackage{listings}
\usepackage{textcomp}
\usepackage{accsupp}
\usepackage{attachfile}
\usepackage{verbatim}
\usepackage[misc]{ifsym}
\usepackage[super]{nth}
\usepackage{dirtree}
\usepackage{booktabs}

% color def
\usepackage{color}
\definecolor{darkred}{rgb}{0.6,0.0,0.0}
\definecolor{darkgreen}{rgb}{0,0.50,0}
\definecolor{lightblue}{rgb}{0.0,0.42,0.91}
\definecolor{orange}{rgb}{0.99,0.48,0.13}
\definecolor{grass}{rgb}{0.18,0.80,0.18}
\definecolor{pink}{rgb}{0.97,0.15,0.45}

\usepackage{tikz}
\usetikzlibrary{shapes.multipart}
\usetikzlibrary {shapes.misc} 
\usetikzlibrary{shadows.blur}
\usetikzlibrary{positioning}
\usetikzlibrary{arrows}
\usetikzlibrary{matrix}
\tikzset{
stack_item/.style={rounded rectangle, rounded rectangle arc length=90,draw, anchor=center, fill=white, minimum width=1cm},
bar/.style={rounded rectangle, minimum width=1.5cm, fill=black, blur shadow={shadow blur steps=5,shadow blur extra rounding=1.3pt}},
message/.style={rectangle, fill=white, color=white, text=black},
bowling_pin/.style={circle, draw, line width=0.6mm, color=red, text=black, minimum width=8mm},
empty/.style={circle, draw, line width=0.6mm, color=white, text=black, minimum width=8mm},
table_cell/.style={rectangle, draw, color=black, text=black, minimum size=12mm}
}

% listings
\usepackage{listings}

% General Setting of listings
\lstset{
	language=Python,
	tabsize=4,
	aboveskip=1em,
	breaklines=true,
	abovecaptionskip=-6pt,
	captionpos=b,
	escapeinside={\%*}{*)},
	frame=single,
	numbers=left,
	numbersep=8pt,
	numberstyle=\tiny,
	morekeywords=[1]{,as,assert,nonlocal,with,yield,self,True,False,None,} % Python builtin
	morekeywords=[2]{,__init__,__add__,__mul__,__div__,__sub__,__call__,__getitem__,__setitem__,__eq__,__ne__,__nonzero__,__rmul__,__radd__,__repr__,__str__,__get__,__truediv__,__pow__,__name__,__future__,__all__,}, % magic methods
	morekeywords=[3]{,object,type,isinstance,copy,deepcopy,zip,enumerate,reversed,list,set,len,dict,tuple,range,xrange,execfile,real,imag,reduce,str,repr,}, % common functions
	morekeywords=[4]{,Exception,NameError,IndexError,SyntaxError,TypeError,ValueError,OverflowError,ZeroDivisionError,}, % errors
	morekeywords=[5]{,ode,fsolve,sqrt,exp,sin,cos,arctan,arctan2,arccos,pi, array,norm,solve,dot,arange,isscalar,max,sum,flatten,shape,reshape,find,any,all,abs,plot,linspace,legend,quad,polyval,polyfit,hstack,concatenate,vstack,column_stack,empty,zeros,ones,rand,vander,grid,pcolor,eig,eigs,eigvals,svd,qr,tan,det,logspace,roll,min,mean,cumsum,cumprod,diff,vectorize,lstsq,cla,eye,xlabel,ylabel,squeeze,}, % numpy / math
	deletekeywords=[2]{next,set},
	deletekeywords=[3]{next,set},
	basicstyle=\ttfamily,
	backgroundcolor=\color{white},
	commentstyle=\color{darkgreen},
	keywordstyle=\color{blue}\bfseries\itshape,
	keywordstyle=[2]\color{blue}\bfseries,
	keywordstyle=[3]\color{grass},
	keywordstyle=[4]\color{red},
	keywordstyle=[5]\color{orange},
	stringstyle=\color{darkred},
	emphstyle=\color{pink}\underbar,
	morestring=[d]{"},
	showstringspaces=false,
	upquote=true,
}


\makeatletter
\lst@RequireAspects{writefile}

% Use \attachfile command to add listing content as paperclip.
\lstnewenvironment{attachedlisting}[1][]{%
	\lst@BeginAlsoWriteFile{#1.txt}%
}
{%
	\lst@EndWriteFile
	\marginpar{\attachfile[appearance=false,icon=Paperclip,mimetype=text/tex]{#1.txt}}
}

% Use accsupp package to add listing content as copyable text.
\lstnewenvironment{copyablelisting}{%
	\lst@BeginAlsoWriteFile{\jobname.txt}%
}
{%
	\lst@EndWriteFile
	\let\verbatim@processline\add@lstline
	\global\let\lstfile\empty
	\verbatiminput{\jobname.txt}%
	\marginpar{(\BeginAccSupp{method=escape,ActualText={\lstfile}}\PaperPortrait\EndAccSupp{})}
}
\def\add@lstline
{\xdef\lstfile{\unexpanded\expandafter{\lstfile}\the\verbatim@line\string^^J}}
\makeatother

\usepackage{titling}
\setlength{\droptitle}{-8em}   % This is your set screw

\title{Postfix Calculator}
\date{}

\begin{document}
\maketitle

\section{Introducción}

En este proyecto implementarás una calculadora para expresiones en notación postfix usando una \texttt{Pila}. El primer paso será calcular el resultado de una expresión postfix. Luego implementaremos un convertidor que tomará una expresión en notación infix y la convertirá a una postfix, para poder usar la calculadora.

\subsection{¿Qué es Postfijo?}
\textit{postfix}, o Notación Polaca Inversa, es una notación matemática donde los operadores se escriben después de sus operandos. A diferencia de la notación \textit{infix} que normalmente usamos para escribir operaciones matemáticas. Una ventaja del postfijo es que no es necesario usar paréntesis para denotar la precedencia de los operadores.

Si tomamos la siguiente expresión \texttt{infix}:

\begin{center}
\texttt{2 + ( 4 + 5 ) / 2 - 8}
\end{center}

Podemos reescribirla en postfix de la siguiente manera:

\begin{center}
\texttt{2 4 5 + 2 / + 8 -}
\end{center}

Nota que los paréntesis ya no están en la expresión, y la precedencia se aplica leyendo los operandos de izquierda a derecha.

\vspace{1cm}
Para obtener el resultado usamos un algoritmo simple utilizando una Pila:
\begin{itemize}
\item Lee cada elemento en la expresión \textit{postfix} de izquierda a derecha.
    \begin{itemize}
    \item Si es un número, lo empujamos a la pila.
    \item Si encontramos un operador, sacaremos 2 números de la pila, aplicaremos el operador a ellos y empujaremos el resultado a la pila. Cuidado con el orden de los operandos, el operando más a la derecha se aplica al operando más a la izquierda.
        \begin{itemize}
        \item[$\circ$] La expresión \texttt{2 4 -} es lo mismo que \texttt{2 - 4} y debería resultar en \texttt{-2}
        \item[$\circ$] La expresión \texttt{9 3 /} es lo mismo que \texttt{9 / 3} y debería resultar en \texttt{3}
        \end{itemize}
    \end{itemize}
\item Al final, nos quedamos con el resultado en la pila.
\end{itemize}

\newpage

\section{Calculadora Postfix}

\subsection{Preparando la Entrada \lstinline|String|}

Implementa una función, \lstinline|preparar_expresion_postfix(string: str) -> list|, que tome un string como parámetro y devuelva una lista de strings. La cadena de entrada será una expresión postfix donde los elementos estarán separados por espacios. Queremos obtener cada elemento individual.
\begin{lstlisting}
elementos = preparar_expresion_postfix("2 4 5 + 2 / + 8 -")
print(elementos)
# Salida:
# ['2', '4', '5', '+', '2', '/', '+', '8', '-']
\end{lstlisting}

\subsection{Obteniendo el Resultado}

Implementa el algoritmo descrito anteriormente para calcular el resultado de una expresión postfix. Solo necesitas tener en cuenta la suma, resta, multiplicación y división, y la entrada siempre contendrá positivos. Escribe una función, \lstinline|evaluar_postfix(elementos: lst) -> float|, que tome la salida de la función anterior y devuelva el resultado de evaluar la expresión.


Aquí hay algunos ejemplos de cómo debería comportarse la función:

\begin{lstlisting}
elementos = preparar_expresion_postfix("2 4 -")
resultado = evaluar_postfix(elementos)
print(resultado)
# salida: -2.0

elementos = preparar_expresion_postfix("9 3 /")
resultado = evaluar_postfix(elementos)
print(resultado)
# salida: 3.0

elementos = preparar_expresion_postfix("2 4 5 + 2 / + 8 -")
resultado = evaluar_postfix(elementos)
print(resultado)
# salida: -1.5

elementos = preparar_expresion_postfix("5 1 - 4 / 2  3 * +")
resultado = evaluar_postfix(elementos)
print(resultado)
# salida: 7.0

elementos = preparar_expresion_postfix("25 2 8 3 3 + - * /")
resultado = evaluate_postfix(elementos)
print(resultado)
# salida: 6.25
\end{lstlisting}

\newpage

\section{Convirtiendo Infix a Postfix}

Las expresiones infix son fácilmente legibles y resolubles por los humanos. Sin embargo, los ordenadores tienen dificultades para diferenciar entre operadores y paréntesis, a menudo requiriendo múltiples pasadas para resolver la operación. Por lo tanto, es más eficiente convertir la expresión a forma postfix antes de la evaluación.

Ahora que tenemos una calculadora funcional, automatizaremos la conversión de expresiones. Terminaremos con un programa que pueda tomar una expresión infix regular y mostrar su resultado.

\subsection{Preparando la Cadena Infix}
El primer paso en la conversión es tokenizar la cadena de entrada. Crea una función, \lstinline|preparar_expresion_infix(string: str) -> list|, que tomando una expresión infix devuelva una lista con los elementos individuales que la componen. Ten en cuenta que la expresión infix puede o no tener espacios separando los elementos.

\begin{lstlisting}
elementos = preparar_expresion_infix("2+(4+5)/2-8")
print(elementos)
# salida: ["2", "+", "(", "4", "+", "5", ")", "/", "2", "-", "8"]

elementos = preparar_expresion_infix("22+897/3")
print(elementos)
# salida: ["22", "+", "897", "/", "3"]

elementos = preparar_expresion_infix("12 * ( 78.9 + 2 )")
print(elementos)
# salida: ["12", "*", "(", "78.9", "+", "2", ")"]
\end{lstlisting}

\subsection{Convirtiendo a Postfix}

Este siguiente paso es el más complejo. Tendremos que tomar la lista de tokens creada por la función anterior y procesarla para obtener la expresión postfix correcta. Necesitaremos usar una Pila y tener en cuenta la precedencia de los operadores mostrada en la Tabla~\ref{tab:operator_precedence}. Escanea la lista de tokens de izquierda a derecha y sigue estos pasos para convertir la expresión:

\begin{enumerate}
\item Si el carácter es un número, ponlo en la expresión de salida.
\item Si el carácter es un operador, mira la pila:
    \begin{itemize}
    \item Si la pila está vacía, el valor de arriba es ``('' o el operador en la parte superior de la pila tiene una precedencia menor que el operador actual, empuja el operador a la pila.
    \item De lo contrario, elimina operadores de la pila mientras tengan más o igual precedencia que el actual, agregando los operadores a la expresión de salida. Finalmente, empuja el operador actual a la pila.
    \end{itemize}
\item Si el carácter es un ``('' empújalo a la pila.
\item Si el carácter es un ``)'' saca elementos de la pila, añadiéndolos a la expresión de salida, hasta que se salga un ``(''. Luego descarta ambos paréntesis.
\item Repite los pasos 1-4 hasta que la lista de tokens esté completamente escaneada.
\item Saca los operadores restantes de la pila y añádelos a la expresión de salida.
\end{enumerate}

Define una función, \lstinline|convertir_a_postfix(infix: str) -> str|, que tome una expresión infix como argumento y devuelva la expresión postfix correspondiente. Tendrás que tokenizar la cadena y aplicar el algoritmo descrito. Recuerda añadir un espacio separador cada vez que añadas algo a la expresión de salida.

\begin{table}[h!tbp]
\centering
\renewcommand{\arraystretch}{1.5}
\begin{tabular}{cc}
\toprule
\textbf{Operador} & \textbf{Precedencia} \\ 
\midrule
$\wedge$ & 3 \\ 
$*\ /$ & 2 \\ 
$+\ -$ & 1 \\ 
\bottomrule
\end{tabular}
\caption{Precedencia de operadores para la conversión de infix a postfix.}
\label{tab:operator_precedence}
\end{table}

\begin{lstlisting}
postfix = convertir_a_postfix("2+(4+5)/2-8")
print(postfix)
# salida: "2 4 5 + 2 / + 8 -"

postfix = convertir_a_postfix("22+897/3")
print(postfix)
# salida: "22 897 3 / +"

postfix = convertir_a_postfix("12 * ( 78.9 + 2 )")
print(postfix)
# salida: "12 78.9 2 + *"
\end{lstlisting}

\section{Crear un Programa}
Desarrolla un programa que solicite al usuario que ingrese una expresión infix y luego muestre el resultado de la expresión, creando una calculadora. El programa pedirá al usuario que ingrese nuevas expresiones hasta que se ingrese la cadena \lstinline|"exit"|. Asegúrate de usar la función \lstinline|main()| como el punto de entrada del programa.

\section{Adiciones al Programa}

\begin{itemize}
\item El programa actualmente admite suma, resta, multiplicación y división. ¿Puedes agregar soporte para la potenciación?
\item Tal y como está el programa, no podemos ingresar valores negativos. Modifica el código para aceptar correctamente cualquier valor negativo ingresado por el usuario.
\item Todas las operaciones en el programa son operaciones binarias\footnote{\href{https://en.wikipedia.org/wiki/Binary_operation}{https://en.wikipedia.org/wiki/Binary\_operation}}, toman dos argumentos para producir la salida. Crea una operación ternaria\footnote{\href{https://en.wikipedia.org/wiki/Ternary_operation}{https://en.wikipedia.org/wiki/Ternary\_operation}} y agrega soporte para ella en el programa. Puedes decidir qué hace y su precedencia.
\end{itemize}


\end{document}